\documentclass[11pt]{article}

\usepackage[utf8]{inputenc}
\usepackage[T1]{fontenc}
\usepackage[a4paper,margin=1in]{geometry}
\usepackage{amsmath}
\usepackage{amssymb}
\usepackage{xcolor}
\usepackage{framed}
\usepackage{enumitem}
\usepackage{hyperref}
\hypersetup{colorlinks=true, linkcolor=blue, urlcolor=blue}

% ---- formatting helpers --------------------------------------------------
\definecolor{cmtbg}{RGB}{244,244,248}
\definecolor{cmtrule}{RGB}{120,120,160}
\definecolor{exrule}{RGB}{40,120,80}

% Reviewer comment: shaded, left-ruled, italic block
\newenvironment{comment}
  {\par\medskip\noindent
   \colorlet{shadecolor}{cmtbg}%
   \def\FrameCommand{\color{cmtrule}\vrule width 2pt\hspace{8pt}}%
   \begin{shaded}\itshape\small}
  {\end{shaded}\par\smallskip}

% Our response
\newcommand{\response}{\par\smallskip\noindent\textbf{Response.}\ }

% Quoted revised manuscript text
\newenvironment{revised}
  {\par\smallskip\noindent
   \def\FrameCommand{\color{exrule}\vrule width 2pt\hspace{8pt}}%
   \MakeFramed{\advance\hsize-\width \FrameRestore}%
   \small\textcolor{exrule}{\textbf{Revised text:}}\ \itshape}
  {\endMakeFramed\par\smallskip}

\newcommand{\rc}[1]{\medskip\noindent\textbf{#1}\par}

\title{\vspace{-1.5cm}\textbf{Response to Reviewers}}
\author{}
\date{}

\begin{document}
\maketitle
\vspace{-2.2cm}

\noindent\textbf{Manuscript:} XHBot: eXplainable Heterophily-aware Graph Neural Networks for Social Bot Detection\\
\textbf{Journal:} EAI Endorsed Transactions on AI and Robotics\\
\textbf{Decision:} Revisions Required

\bigskip

\noindent Dear Associate Editor and Reviewers,

\medskip

We thank the Associate Editor and both reviewers for their careful and constructive
evaluation of our manuscript. The comments have substantially improved the clarity,
rigour, and presentation of the work. We have revised the manuscript thoroughly and
respond to every point below. Throughout this letter we use the following convention:
reviewer comments appear in \emph{shaded italic boxes}, our responses follow in
regular text, and the most relevant added or modified manuscript text is reproduced
in \textcolor{exrule}{green-ruled boxes}. All changes compile cleanly; the revised
LaTeX source and the compiled PDF accompany this letter.

\section*{Summary of the principal changes}
\begin{enumerate}[leftmargin=*,itemsep=1pt]
  \item Removed the duplicated Conclusion section and the inconsistent ``+27.37\%''
        figure; the manuscript now reports a single, consistent set of improvements
        (9.64\% on TwiBot-20, 13.15\% on TwiBot-22).
  \item Rewrote the abstract as a single flowing paragraph without structured labels.
  \item Added a spectral-theory justification for SGTR based on the graph Dirichlet
        energy, and an explicit comparison with CARE-GNN and RABot.
  \item Rebuilt the ablation study to ablate each of the three THCA channels
        independently and to include \emph{w/o CPD} and \emph{w/o SGTR} variants;
        corrected the erroneous ``Dual-Channel'' terminology.
  \item Added a baseline-fairness statement (identical splits, identical imbalance
        sampling across all models, re-tuned hyperparameters).
  \item Added a qualitative forensic case study for the HFE module.
  \item Replaced the cluttered architecture figure with a clean vector diagram that
        explicitly shows the three THCA channels; clarified the robustness figure and
        expanded its analysis.
  \item Softened absolute/exaggerated language throughout; standardised the
        bibliography (DOIs, author lists) and in-text citation style; and specified
        the previously missing hyperparameters.
\end{enumerate}

% =========================================================================
\section*{Response to Reviewer 1}

\rc{Comment 1 --- Novelty of SGTR relative to CARE-GNN and RABot.}
\begin{comment}
One of the core modules, ``Spectral-Guided Topology Refinement (SGTR),'' aims to
prune camouflage edges by calculating edge spectral energy. However, this concept is
highly similar in both motivation and implementation to existing works, such as the
label-aware similarity metric in CARE-GNN and the reinforcement learning-guided edge
filtering in RABot. \ldots Is merely replacing the similarity metric with ``spectral
energy'' sufficient to constitute a core innovation?
\end{comment}
\response We agree the original text did not adequately differentiate SGTR. We have
(i) added a spectral-theory derivation grounding the ``spectral'' terminology (see
Reviewer~2, Comment~3), and (ii) added a dedicated paragraph, \emph{``Relation to
prior edge-refinement methods,''} in Section~3.2, and a contrasting sentence in the
related-work RABot paragraph. We now position SGTR as a label-free, theoretically
motivated, single-step differentiable alternative, and we measure its isolated
contribution in the ablation (Table~2, \emph{w/o SGTR}, $-0.0494$~F1).
\begin{revised}
SGTR is related to, but distinct from, two lines of prior work. CARE-GNN identifies
informative neighbours using a \emph{label-aware} similarity measure trained with
node labels and a reinforcement-learning neighbour selector; RABot discards edges
below a reliability threshold that is itself tuned online by a separate
reinforcement-learning agent. In contrast, SGTR (i) is \emph{label-free}, scoring
edges purely from feature discrepancy and thus avoiding the propagation of label
noise into the topology; (ii) is grounded in the graph Dirichlet-energy
decomposition, giving the score a spectral interpretation rather than a heuristic
similarity; and (iii) is a single differentiable pre-aggregation step trained
end-to-end with the detector, rather than a separately optimised RL policy.
\end{revised}

\rc{Comment 2 --- Baseline fairness and imbalance sampling.}
\begin{comment}
\ldots the paper mentions that XHBot uses ``Imbalance Sampling'' to handle class
imbalance. The reviewer questions whether this strategy was applied to all baseline
models as well. If the baseline models were not evaluated using the same data
sampling strategy, the performance improvements might be partially attributed to data
processing techniques rather than the inherent superiority of the model architecture.
\end{comment}
\response This is an important point. We added Section~4.1.4, ``Baseline
Implementation and Fair Comparison,'' which states that imbalance-aware sampling was
applied uniformly to all baselines that support it, that all models share identical
splits/features/graph, and that baseline hyperparameters were re-tuned per dataset.
The \emph{w/o Imbalance Sampling} ablation row isolates this component for XHBot.
\begin{revised}
\ldots the class-imbalance handling used by XHBot (imbalance-aware sampling) was
applied uniformly to all baselines that support it, so no model receives an advantage
from sampling alone\ldots For each baseline we started from the authors' reference
implementation and hyperparameters where available, and additionally re-tuned the key
hyperparameters \ldots on the validation set of each dataset. Consequently, the
comparison reflects re-optimised baselines on the datasets used in this study rather
than figures transcribed from their original papers.
\end{revised}

\rc{Comment 3 --- Ablation study flaws (``Dual-Channel'' vs.\ Tri-channel; missing CPD).}
\begin{comment}
The ablation study (Table 2) has obvious flaws. For example, it removes ``Dual-Channel
Aggregation,'' whereas the THCA module is actually ``Tri-channel.'' \ldots Furthermore,
the paper lacks an ablation study for the ``Contrastive Prototype Disentanglement
(CPD)'' module \ldots Table 2 does not include a ``w/o CPD'' variant.
\end{comment}
\response We rebuilt Table~2 entirely. It now ablates each THCA channel independently
(\emph{w/o Homophilic}, \emph{w/o Heterophilic}, \emph{w/o Self-Identity}), retains a
homophilic-only extreme for reference, and adds \emph{w/o SGTR} and \emph{w/o CPD},
with a new $\Delta$F1 column. The erroneous ``Dual-Channel'' label has been removed
throughout. Removing the heterophilic channel produces the largest single-channel
drop ($-0.1244$~F1); \emph{w/o CPD} quantifies the disentanglement contribution
($-0.0424$~F1), showing CPD improves separability beyond THCA alone.

\rc{Comment 4 --- HFE explainability is evaluated only quantitatively.}
\begin{comment}
\ldots metrics like Fidelity and Sparsity \ldots are far from sufficient for a
human-oriented task like ``forensic-level explanation.''
\end{comment}
\response We added Section~4.6, ``Qualitative Forensic Case Study,'' which walks a
representative flagged account through the three HFE evidence levels (instance
subgraph, community ring, diagnostic motif) with a human-readable interpretation, and
we explicitly acknowledge that automated metrics are necessary but not sufficient.
\begin{revised}
\ldots these three levels turn an opaque probability into a reviewable narrative
(``this account, together with these 11 others, follows many unrelated humans in a
coordinated star pattern''), which is the form of evidence a moderation team can act
on and document. We emphasise that this case study is illustrative; a rigorous
assessment of whether such explanations improve human moderator decisions would
require a dedicated user study, which we identify as future work.
\end{revised}

\rc{Comment 5 --- Inconsistent ``+27.37\%'' claim in the conclusion.}
\begin{comment}
In the conclusion (Section 7), the paper claims ``securing a massive +27.37\% F1 score
improvement.'' This figure is completely inconsistent with the ``+9.64\%'' and
``+13.15\%'' reported in the main experimental section (Section 4.2).
\end{comment}
\response The inconsistency arose from an accidentally duplicated Conclusion section.
We deleted the duplicate section and the erroneous figure. The manuscript now reports
only the values supported by Section~4.2 (9.64\% on TwiBot-20, 13.15\% on TwiBot-22),
stated consistently in the abstract, introduction, results, and conclusion.

\rc{Comment 6 --- Figure clarity (architecture and robustness).}
\begin{comment}
As the model architecture diagram, Fig. 1 is highly informative but somewhat cluttered
\ldots the internal details of certain modules (e.g., the three channels of THCA) are
not clearly expressed. The title of Fig. 4 is ``Model robustness \ldots'' but the
specific meanings of the curves \ldots are unclear in the legend, and the analysis of
this figure in the main text is rather brief.
\end{comment}
\response Figure~1 was redrawn as a clean vector (TikZ) diagram with an explicit
left-to-right data flow and a dedicated, labelled THCA block showing all three
channels (homophilic / heterophilic / self-identity). Figure~4 (robustness) was
regenerated with a titled legend, an explicit x-axis definition of ``camouflage
degree,'' distinct markers, and end-point value annotations; its analysis was expanded
from a single sentence to a full paragraph quantifying each model's degradation.

\rc{Comment 7 --- Exaggerated language.}
\begin{comment}
The paper frequently uses absolute terms such as ``fundamentally insufficient'' and
``impenetrable black boxes'' \ldots it is recommended to use more objective and neutral
language.
\end{comment}
\response We revised absolute phrasing throughout, including ``impenetrable black
boxes'' $\rightarrow$ ``provide limited rationale,'' ``fundamentally insufficient''
$\rightarrow$ ``often insufficient,'' ``provably insufficient'' $\rightarrow$ ``often
insufficient,'' and ``comprehensively dominates / massive margin / weaponizes''
$\rightarrow$ neutral, quantitative descriptions.

\rc{Comment 8 --- Reference inconsistencies (DOIs, author lists).}
\begin{comment}
In the reference list, some entries include DOI links while others do not. The
formatting of the author lists for some conference papers is also inconsistent.
\end{comment}
\response We standardised the bibliography. The arXiv entries now use a uniform
\texttt{journal = \{arXiv preprint\}} + \texttt{doi} format with redundant
URL/identifier fields removed; a DOI was added to the GNNExplainer entry; and author
lists use a uniform \texttt{Last, First} format. The two genuine no-DOI NeurIPS
proceedings entries retain a canonical URL. We also fixed a BibTeX warning by removing
a conflicting \texttt{number} field in the RGT entry.

\rc{Comment 9 --- Optimise figures and language.}
\response Addressed jointly with Comments~6 and~7 above.

% =========================================================================
\section*{Response to Reviewer 2}

\subsection*{Major comments}

\rc{Major 1 --- Abstract format.}
\begin{comment}
The abstract uses explicit labels like ``INTRODUCTION'' and ``METHODS.'' This is
non-standard for most journals. It should be rewritten as a single, flowing paragraph.
\end{comment}
\response The abstract has been rewritten as a single flowing paragraph; all
structured labels have been removed.

\rc{Major 2 --- Structure and redundancy (two Conclusions; verbose intro/related work).}
\begin{comment}
\ldots there are two separate ``Conclusion'' sections (Section 6 and Section 7) that
repeat similar points. The authors must consolidate these \ldots the introduction and
related work sections are overly verbose and could be significantly streamlined.
\end{comment}
\response The duplicated Conclusion was removed, leaving a single cohesive conclusion.
We condensed the statistics-heavy opening of the introduction and substantially
shortened the related-work section (notably the RABot and FEDRIO paragraphs) while
preserving the technical contrasts.

\rc{Major 3 --- Clarity of the SGTR module; ``Spectral-Guided'' is misleading.}
\begin{comment}
\ldots The text claims that cross-class connections manifest as ``high-frequency
noise,'' yet Equation (1) defines a ``spectral energy'' \ldots using a learnable neural
network on node features, not a spectral graph filter. The term ``Spectral-Guided'' is
therefore misleading. \ldots How does learning an edge weight based on feature
discrepancy relate to spectral graph theory?
\end{comment}
\response We added a ``Spectral motivation'' paragraph that makes the connection
explicit. The graph Dirichlet energy (Laplacian quadratic form) decomposes over edges,
so the squared feature discrepancy used in the score is exactly the per-edge
contribution to the high-frequency (Dirichlet) energy; down-weighting high-energy
edges therefore attenuates the high-frequency band that drives over-smoothing. The
learnable projection is presented as a \emph{calibrated, data-dependent estimate} of
this quantity, which justifies the terminology.
\begin{revised}
For a graph signal $\mathbf{f}\in\mathbb{R}^N$ and the combinatorial Laplacian
$L=D-A$, the Dirichlet energy admits the edge decomposition
$\mathbf{f}^\top L \mathbf{f} = \sum_{(u,v)\in\mathcal{E}}(f_u-f_v)^2$, which measures
the high-frequency content of $\mathbf{f}$. \ldots the squared feature discrepancy
$(\mathbf{x}_u-\mathbf{x}_v)^2$ is the per-edge contribution to the total feature
Dirichlet energy. \ldots SGTR scores edges by this quantity and suppresses the
highest-energy ones, which is equivalent to attenuating the high-frequency band that
drives over-smoothing in low-pass GCN aggregation.
\end{revised}

\rc{Major 4 --- Justification for the CPD module (add ``w/o CPD'').}
\begin{comment}
\ldots The ablation study in Table 2 \ldots does not include a variant ``w/o CPD.'' The
authors must add this ablation \ldots Without it, it is unclear if the performance gains
are driven by the THCA module alone.
\end{comment}
\response Added (see Reviewer~1, Comment~3). Table~2 now includes \emph{w/o CPD}
($-0.0424$~F1), demonstrating that CPD contributes beyond THCA alone, primarily by
sharpening class separability.

\rc{Major 5 --- Explainability evaluation: add qualitative case studies.}
\begin{comment}
\ldots The authors should include qualitative case studies. For example, show a specific
bot account flagged by the model, present the subgraph identified by HFE as the
``cause,'' and provide a human-interpretable explanation of why that subgraph is
suspicious \ldots
\end{comment}
\response Added (see Reviewer~1, Comment~4). Section~4.6 presents a concrete flagged
account, the HFE subgraph as the cause, and a human-interpretable explanation of the
``star-formation'' botnet pattern, exactly as suggested.

\rc{Major 6 --- Baseline selection and fairness.}
\begin{comment}
\ldots It would strengthen the paper to discuss whether the hyperparameters for these
baseline models were re-optimized on the specific datasets used in this study, or if
the results were taken from their original papers. Ensuring a fair comparison is
critical.
\end{comment}
\response Addressed in Section~4.1.4 (see Reviewer~1, Comment~2): baseline
hyperparameters were re-tuned on this study's datasets rather than transcribed from
the original papers, and the full comparison protocol is now stated explicitly.

\subsection*{Minor comments}

\rc{Minor 1 --- Figure 1 density.}
\response The architecture diagram was redrawn as a clean vector figure with a clearer
data flow and an explicit THCA grouping (see Reviewer~1, Comment~6).

\rc{Minor 2 --- Missing hyperparameters (hidden dimension, number of layers).}
\begin{comment}
\ldots does not specify the hidden dimension size or the number of layers for the GNN
backbone. These details are necessary for reproducibility.
\end{comment}
\response Section~4.1.3 now specifies the number of THCA layers ($L=2$), hidden
dimension ($d_h=64$), dropout ($0.5$), CPD sub-space dimension ($d'=32$), InfoNCE
temperature ($\tau_{nce}=0.1$), loss weights ($\lambda_{MI}=0.1$,
$\lambda_{NCE}=0.5$), and the SGTR threshold ($\tau=0.5$).

\rc{Minor 3 --- Confusing notation in Equation (2).}
\begin{comment}
Typo in Equation (2): the notation is slightly confusing. Please clarify.
\end{comment}
\response The soft-pruning rule was rewritten in explicit piecewise (cases) form,
making the thresholding behaviour unambiguous: low-energy edges keep their original
weight, while high-energy edges are attenuated by a factor $(1-E)$.
\begin{revised}
$\tilde{A}^{(r)}_{u,v} = A^{(r)}_{u,v}\,(1-E^{(r)}(u,v))$ if $E^{(r)}(u,v)>\tau$, and
$\tilde{A}^{(r)}_{u,v} = A^{(r)}_{u,v}$ otherwise.
\end{revised}

\rc{Minor 4 --- Citation formatting.}
\begin{comment}
There are some inconsistencies in the citation style (e.g., ``[17] proposes MGDIL'' vs
``The work of [17] proposes\ldots''). Please standardize the citation format throughout.
\end{comment}
\response We standardised in-text citations to the ``Author et al.~[n]'' form when a
work is the subject of a sentence (e.g., ``Qiao et al.~[17] propose MGDIL,'' ``Fu et
al.~[8],'' ``He et al.~[9],'' ``Dang and Nguyen~[15]''), removing the inconsistent
bare-number and ``The work of \ldots'' constructions.

\bigskip
\noindent We believe these revisions comprehensively address the reviewers' concerns
and meaningfully strengthen the manuscript. We thank the Associate Editor and both
reviewers again for their time and valuable feedback.

\medskip
\noindent Sincerely,\\
The Authors

\end{document}
